\section{Discussion}
\label{sec:discussion}

\paragraph{What the results mean.} In this formulation the current LSP
assignment is genuine state, and a move constrains the future through
hold-downs, reversal costs and coupled protected-class legality. Yet a learner
that ignores the future reaches the return of a per-interval optimizer, and a
learner trained with an effective horizon of about 200 intervals
($1/(1-\gamma)$, 17 hours) does worse than both. The closed-loop oracle ladder
shows why there is little to gain from looking ahead: even with perfect
foresight, a controller that re-plans every interval gains only a few points
from lookahead beyond the next interval. When a move's effect is felt in the
interval in which it is made, the immediate reward ranks moves almost as a
longer criterion would; what is not myopic---anticipating bursts, failures
and diurnal change---is not provided by an observation without a clock or
forecasts. A sequential learner pays the price of a long-horizon objective
(noisy advantages, unstable optimization) without access to the signal that
would make it worthwhile.

\paragraph{Long horizons can hide small immediate costs.} The concrete failure
on three of five PPO roots is instructive for RL practice. Reconfiguration costs are small, immediate and
certain; congestion rewards are large, delayed and noisy. With $\gamma=0.995$
PPO's advantage estimates are dominated by the latter and fail to attribute
the former, so the policy settles into flipping a demand whenever the
hold-down allows. We have not measured the advantage noise directly. The
evidence for this explanation is that setting $\gamma=0$ removes the
oscillation and closes the gap; against it, the critic is not poorly fitted
in aggregate (explained variance at the last update
\PPOEVDefaultMin--\PPOEVDefaultMax{} with $\gamma=0.995$, versus
\PPOEVZeroMin--\PPOEVZeroMax{} with $\gamma=0$), so the failure would have to
lie in the small residual differences between moves rather than in the value
scale; and the PPO roots that do not oscillate fall short through lower network utility
rather than churn. The discount factor is thus a decisive hyper-parameter here, and the
default $\gamma\approx0.99$ that most RL-for-networking work inherits is not a
neutral choice when costs and benefits live on different time scales.

\paragraph{When is sequential RL justified?} Our evidence points to three
conditions, each testable before choosing an algorithm: (i) actions whose
effects are delayed relative to the decision, so that the immediate reward
does not contain them (\S\ref{sec:rq5}); (ii) exogenous dynamics that are
predictable from the observation, so that anticipation can be learned; and
(iii) decisions that change the future feasible set substantially. The
closed-loop oracle ladder of \S\ref{sec:sequentiality} bounds the value of
lookahead under perfect foresight and the frozen rollout measures how much
cost amortization changes the ranking of moves, both from a simulator without
learning; neither measures (iii) beyond the bound, because their continuation
never moves again. We recommend reporting both before choosing an algorithm.

\paragraph{When should simpler methods be preferred?} When effects are
immediate, a myopic learner---or, if a model of the objective is available, a
per-interval optimizer---reaches a higher return than PPO here. Inference cost
does not discriminate at a five-minute interval (the bandit needs
\InfBanditMs{}\,ms per decision, PPO \InfPPOMs{}\,ms, MILP-track
\InfMilpMs{}\,ms). The bandit and MILP-track trade differently: MILP-track
attains higher network utility and pays more in move costs, which its
objective does not contain; the bandit reroutes \ChurnRatioBanditMilp{} times
as often as MILP-track (\TestBanditReroutes{} vs.\ \TestMilpReroutes{} per
hour) at a return we cannot distinguish from MILP-track's, which matters
operationally because every LSP move is signalling and risk. A churn-aware
optimizer would be the natural stronger baseline.

\paragraph{Networking implications.} Teal and DOTE treat TE as a one-shot
problem per traffic matrix. Our results support that treatment in a setting
with persistent state, incremental changes and stability timers, where one
might have expected sequential structure to matter---as long as a change takes
effect within the interval in which it is decided, and in a heavily loaded
regime where TE mostly decides which traffic to lose (\S\ref{sec:limitations}). They also show that churn
metrics (reroutes, reversals) must be reported with the objective: PPO's best
root reaches a higher network utility than any bandit root (\PPOUtilMax{}
vs.\ at most \BanditUtilMax) but would be unacceptable to an operator because of its oscillation.

\paragraph{Generalization.} We make no claim beyond this simulator. Larger
topologies, elastic traffic, real traffic matrices and control-plane latency
all plausibly increase temporal coupling; the diagnostic and the
delayed-effect experiment indicate how to test that.
